\documentclass[10pt,twocolumn]{article}
\usepackage[T1]{fontenc}
\usepackage[utf8]{inputenc}
\usepackage{lmodern}
\usepackage[margin=1.8cm,top=2.4cm,bottom=2cm,columnsep=0.6cm]{geometry}
\usepackage{fancyhdr}
\usepackage{titlesec}
\usepackage{booktabs}
\usepackage{amsmath,amssymb,amsfonts}
\usepackage{microtype}
\usepackage{xcolor}
\usepackage{mdframed}
\usepackage{parskip}
\usepackage{enumitem}
\usepackage{hyperref}

\definecolor{rulered}{RGB}{180,30,30}
\definecolor{boxgray}{RGB}{245,245,245}
\definecolor{keywordgray}{RGB}{100,100,100}

\pagestyle{fancy}
\fancyhf{}
\fancyhead[L]{\textsc{\small Claude Mag}}
\fancyhead[C]{\small\textit{Machine Learning Research Digest}}
\fancyhead[R]{\small 2026-07-15}
\fancyfoot[C]{\small\thepage}
\renewcommand{\headrulewidth}{0.8pt}

\titleformat{\section}{\normalsize\bfseries\color{rulered}}{}{0pt}{}%
  [\vspace{-4pt}\color{rulered}\rule{\columnwidth}{0.4pt}]
\titlespacing{\section}{0pt}{8pt}{4pt}

\hypersetup{colorlinks=true,urlcolor=rulered,linkcolor=black}

\begin{document}
\setlength{\parindent}{0pt}
\setlength{\parskip}{4pt}

{\small\color{keywordgray}\textbf{Keywords:}
  efficient inference \textbullet{}
  diffusion language models \textbullet{}
  speculative decoding \textbullet{}
  throughput}\par\vspace{4pt}

{\LARGE\bfseries\raggedright
  Nemotron-Labs-Diffusion}\par\vspace{4pt}

{\large\itshape\raggedright
  NVIDIA Unifies Autoregressive and Diffusion Language Modeling
  Into a Single Self-Speculating Architecture That Eliminates the
  Need for a Separate Draft Model}\par\vspace{6pt}

{\small\textbf{By} Yonggan Fu, Lexington Whalen, Abhinav Garg,
  Chengyue Wu, et al.\ (NVIDIA Research)}\par\vspace{2pt}

{\small\color{keywordgray}arXiv:2607.05722 \textbullet{}
  July 7, 2026}
\par\vspace{2pt}\noindent{\color{rulered}\rule{\columnwidth}{0.6pt}}\vspace{6pt}

Modern speculative decoding relies on a separate draft model to propose
candidate tokens that a larger verifier model then accepts or rejects.
Maintaining two models doubles deployment cost and complicates serving
infrastructure. Nemotron-Labs-Diffusion eliminates the draft model
entirely by training a single set of weights to operate in three modes:
standard autoregressive (AR) decoding, parallel diffusion decoding, and
self-speculation, where the diffusion head drafts while the AR head
verifies---from the same checkpoint, with no weight switching.

\begin{mdframed}[backgroundcolor=boxgray,linecolor=rulered,linewidth=1pt,
    innerleftmargin=6pt,innerrightmargin=6pt,
    innertopmargin=6pt,innerbottommargin=6pt]
\textbf{\small AT A GLANCE}\par\vspace{4pt}
\begin{description}[leftmargin=0pt,labelsep=4pt,itemsep=2pt,parsep=0pt,
    font=\small\bfseries]
  \item[Problem:] Speculative decoding requires a separate draft model,
    doubling parameter budget and serving complexity.
  \item[Why it matters:] Throughput is the primary cost driver for
    large-scale LLM deployment; a 4--6$\times$ improvement directly
    reduces inference cost.
  \item[Method:] Joint AR-diffusion training objective; at inference,
    the diffusion head proposes draft token blocks and the AR head
    verifies, all within one set of weights.
  \item[Result:] 6.82 accepted tokens per step (vs.\ Eagle3's 2.75);
    4$\times$ higher throughput than Qwen3-8B on SPEED-Bench.
  \item[Evidence:] Evaluated on SPEED-Bench, MT-Bench, MMLU, and
    HumanEval across 3B, 8B, and 14B model sizes.
\end{description}
\end{mdframed}

\section{The Big Picture}

The dominant paradigm for efficient large language model inference is
speculative decoding: a small, cheap draft model proposes a block of
candidate tokens, and the large model verifies the entire block in a
single forward pass. When the acceptance rate is high, this delivers
near-linear speedups. The problem is that the draft model is a
separate artifact that must be trained, maintained, and co-deployed
with the verifier, doubling operational complexity.

Multi-token prediction (MTP) sidesteps the separate model by training
the main model to predict multiple future tokens simultaneously using
auxiliary heads. But MTP heads predict tokens autoregressively from
a fixed left-to-right prior, limiting their ability to propose coherent
multi-token continuations.

Nemotron-Labs-Diffusion observes that the diffusion training objective
---which reconstructs any masked subset of tokens conditioned on
unmasked context---is exactly the capability needed for a
high-quality internal drafter. A model trained jointly with AR and
diffusion objectives learns to both predict the next token from the left
and fill any arbitrary masked block, making its diffusion head a natural
draft model and its AR head a natural verifier.

\section{Why Existing Approaches Fall Short}

Eagle3, the leading separate-draft speculative decoding system, achieves
2.75 accepted tokens per step. Its primary weakness is that the draft
model must be trained separately and kept synchronized with the verifier
model through costly co-training or periodic retraining. Every new
verifier checkpoint requires a corresponding draft model update.

MTP-based approaches (as used in DeepSeek-V3) train auxiliary prediction
heads on the verifier model itself but restrict each head to predicting
a single additional token conditioned on prior outputs. The heads predict
tokens serially rather than in parallel, and their training signal is
weaker because they never reconstruct arbitrary masked subsets---they
only predict a fixed number of positions ahead.

Standard diffusion LMs (as in MDLM and Plaid) have high generation
quality in parallel mode but suffer from a quality-speed tradeoff:
fewer diffusion steps produce lower quality, while more steps reduce
the throughput advantage over AR models.

\section{Core Mechanism}

Nemotron-Labs-Diffusion resolves these tensions by exploiting the
complementarity between AR and diffusion objectives. The AR objective
provides strong left-to-right linguistic priors that make verification
reliable; the diffusion objective provides parallel reconstruction
capability that makes drafting both fast and coherent.

In self-speculation mode, the model executes a draft-verify cycle.
The diffusion head, given the current context, denoises a block of
$B$ masked positions in parallel using one or two diffusion steps.
The AR head then scores each proposed token in a single forward pass
(matching the cost of verifying $B$ tokens in standard speculative
decoding). Tokens above the acceptance threshold are committed; the
rest trigger a fresh diffusion draft.

\section{Technical Formulation}

The joint training loss is:
\[
\mathcal{L} = \alpha\,\mathcal{L}_{\text{AR}} +
  (1-\alpha)\,\mathcal{L}_{\text{diff}}
\]
where $\mathcal{L}_{\text{AR}}$ is the standard cross-entropy on
the next token and $\mathcal{L}_{\text{diff}}$ is the masked token
reconstruction loss at variable mask rates
$\rho \sim \text{Beta}(a, b)$:
\[
\mathcal{L}_{\text{diff}} = -\mathbb{E}_{\rho, \mathcal{M}}
  \sum_{i \in \mathcal{M}} \log p_\theta(x_i \mid x_{\backslash\mathcal{M}})
\]
with $\mathcal{M}$ the random mask set of expected size $\rho n$.
At $\alpha = 0.5$, both objectives contribute equally; ablations
confirm this is the Pareto-optimal balance.

Token acceptance in self-speculation mode follows the standard
speculative decoding criterion: draft token $x_d$ is accepted with
probability $\min(1,\, p_{\text{AR}}(x_d) / p_D(x_d))$, preserving
the AR distribution exactly.

\section{Learning or Inference Procedure}

\begin{enumerate}[itemsep=2pt,parsep=0pt]
  \item Train with joint AR+diffusion loss using standard next-token
    and masked-span reconstruction objectives.
  \item Fine-tune an additional lightweight MLP sampler on diffusion
    sampling trajectories to reduce required denoising steps to 1--2.
  \item At inference, select mode based on deployment constraints:
    AR mode for minimum per-token latency; diffusion mode for maximum
    throughput at batch size 1; self-speculation for the best
    throughput-quality tradeoff at any batch size.
  \item Self-speculation loop: draft $B=8$ tokens via 1-step diffusion;
    score with AR head; commit accepted tokens; redraft rejected suffix.
\end{enumerate}

\section{What the Guarantee Says}

The acceptance-based guarantee from standard speculative decoding
carries over: accepted tokens are drawn exactly from the AR model's
distribution, so Nemotron-Labs-Diffusion-self-speculation produces
outputs statistically identical to pure AR sampling. No quality
degradation is introduced by the self-speculation mechanism.

The throughput guarantee is empirical: on SPEED-Bench with SGLang on
NVIDIA GB200, Nemotron-Labs-Diffusion-8B achieves 4$\times$ higher
throughput than Qwen3-8B at equivalent quality (within 1\% on MT-Bench,
MMLU, and HumanEval).

\section{Experimental Findings}

\begin{center}
\small
\begin{tabular}{lcc}
\toprule
System & Accepted toks/step & Throughput \\
\midrule
Qwen3-8B (AR only) & 1.00 & $1.0\times$ \\
DeepSeek-V3 MTP & 2.11 & $1.9\times$ \\
Eagle3 (separate draft) & 2.75 & $2.6\times$ \\
\textbf{Nemotron-8B (ours)} & \textbf{6.82} & \textbf{4.0$\times$} \\
\bottomrule
\end{tabular}
\end{center}

Quality scores (MT-Bench / MMLU / HumanEval) for Nemotron-8B are
within 1\% of Qwen3-8B AR baseline across all benchmarks. The 3B
model variant outperforms the 7B Eagle3 draft+8B verifier pair on
throughput while using 5B fewer total parameters.

\section{Ablations and Interpretation}

\textbf{Removing diffusion training} drops accepted tokens per step
from 6.82 to 4.1, confirming the diffusion head's centrality to
draft quality. This is the most significant ablation result.

\textbf{$\alpha$ sweep:} Higher $\alpha$ (more AR weight) improves
text quality metrics by up to 0.3\% but reduces acceptance rate by
up to 1.4 tokens per step. $\alpha = 0.5$ is the Pareto-optimal
balance.

\textbf{Mask rate schedule:} Fixed mask rate $\rho = 0.15$ (like
BERT) reduces accepted tokens per step by 1.2 relative to the
$\text{Beta}(1, 5)$ schedule, which samples mask rates across the
full $[0, 1]$ range and builds more versatile reconstruction capability.

\textbf{Self-speculation vs.\ MTP:} At equal parameter budgets,
self-speculation accepts 6.82 tokens per step versus MTP's 2.11,
because MTP heads can only predict a fixed number of positions while
the diffusion head can fill any masked block with a single forward pass.

\section*{Reference}

Yonggan Fu et al.\
\textbf{Nemotron-Labs-Diffusion: A Tri-Mode Language Model Unifying
Autoregressive, Diffusion, and Self-Speculation Decoding}.
arXiv:2607.05722, July 2026.
\url{https://arxiv.org/abs/2607.05722}

\end{document}
